\section{The Good Shepherd}

The Gospel for the Second Sunday after Easter is John 10:11-15\footnote{\url{http://www.biblestudytools.com/rhe/john/passage.aspx?q=John+10:11-30}}, the parable of the Good Shepherd.

In this passage, Jesus compares himself to the Good Shepherd who watches over his sheep and would give his life for them. The Shepherd is clearly the Logos, the Word made flesh. The flock is not so clear. It is obviously not a mechanical aggregate of individual sheep who just happen to come together. On the contrary, they form a unity, which derives from the primordial unity of the divine Logos.

In the Sophiology developed and refined by \textbf{Vladimir Solovyov}, the flock, then, is the unity produced by the Logos, that is, Sophia, the world soul or ideal humanity. Sophia is the mediatrix between the multiplicity of individuals (the real content of her Life) and the absolute Divine unity (her ideal principle).

Furthermore, Jesus says, ``I am the good Shepherd and I know Mine, and Mine know Me, as the Father knoweth Me, and I know the Father.'' Note in this passage the emphasis is on ``knowing'', not on faith or believing. And not just a common knowing, but analogously to how the Logos knows the Father. This type of knowing is superior to rational thinking, it is, rather, a direct intuition of identity, for the sheep know who they are and whom they belong to. They have made themselves meek in order to be able to follow.

When this knowledge is lost, the sheep separate from the flock and become prey to the wolves. With the forgetfulness of Divine unity, that is, when the hireling is watching over the flock, Solovyov explains

\begin{quotex}
the world soul ceases to unite all within herself, all things lose their common bond and the unity of cosmic creation breaks up into a multitude of separate elements.
\end{quotex}


As to the separation of an individual, he continues:

\begin{quotex}
when it is detached from the whole by exciting in itself its own peculiar will, the particular elements of the universal organism lose their common bond in the world soul and, left to themselves, are doomed to discordant, egoistic existence. The root of such existence is evil and the fruit, suffering.
\end{quotex}



\flrightit{Posted on 2010-04-18 by Cologero}

\begin{center}* * *\end{center}

\begin{footnotesize}
\begin{sffamily}
\texttt{GFJF on 2010-04-19 at 00:42 said:}

This is just what I came to ask you about. Until recently I hadn't payed much attention to the traditional aspects of Christianity, the result being that I'm now playing catch-up. However I digest the Neoplatonic terms quite well, so thank you for that. Since Christ, properly understood, is the Logos, it follows that the Father is the One, and the Holy Spirit is the All-Soul. One question I have is this: in the Lord's Prayer it is said that God is in Heaven; that implies the priority of Heaven over God, does it not? It seems to me that Heaven merges in some way with God-as-One, and Father (as demiurge) merges in some way with Christ-as-Logos.

\hfill

\texttt{Cologero on 2010-04-20 at 00:07 said:}

I know some gnostic sects tried to identify the Holy Spirit with Sophia, or the world soul. But they are really on different orders; the spirit is divine and ideal, the world soul is manifesting the divine idea as the cosmic process. This was briefly addressed in another comment.

The Trinity is more complex and will have to be dealt with separately.

\hfill

\texttt{GFJF on 2010-04-20 at 11:39 said:}

Ah. I see I'm dealing with these ideas in the wrong way, treating them primarily as words.

\end{sffamily}
\end{footnotesize}

